\section{Security, Privacy and Ethical Considerations}
\label{sec:security}

\paragraph{Threat model} We consider four threats: (i)~instructions embedded in page content that try to redirect the agent (indirect prompt injection~\cite{greshake2023indirect}); (ii)~high-impact actions taken on the user's behalf without consent, such as posting, sending or purchasing; (iii)~damage to the local system once file and process actions exist; and (iv)~disclosure of user data to parties the user did not intend. The user and the local machine are trusted; websites are not. We do not consider a compromised host or a malicious local model.

\paragraph{Mitigations in the current code} Element text, accessible labels and placeholders pass through a sanitizer that replaces 17 patterns of instruction-override, delimiter, exfiltration and role-manipulation phrasing with a marker before the planner sees them; retrieved knowledge chunks are filtered the same way and wrapped in an untrusted-content banner. Before any navigation, a risk gate pauses tasks whose parsed action is posting, sending e-mail, purchasing, deleting or transferring, or whose risk level is high or critical, and records an approval request; the gate is checked again when completion is verified. The service binds to the loopback interface and accepts browser origins from localhost only. Every model call is logged with its destination, a locality flag and request and response sizes. Observability events are scrubbed of credential-like keys and values before they are broadcast.

\paragraph{What these mitigations do not cover} The sanitizer is a denylist and therefore bypassable by paraphrase, other languages or encodings~\cite{liu2024formalizing}; the page title, URL, extracted text and screenshots reach models unfiltered. The risk level that triggers approval is produced by the same model that reads untrusted content, so an injection that lowers it would also bypass the gate; approval is per task, and once granted it covers every later action. The audit log records model calls only; browser traffic, which necessarily carries whatever the user asks the agent to type into a site, is not logged, and the configured model address is recorded as local or remote but not restricted to the local host. A privacy statement for this system should therefore read: \emph{by default, model inference runs on the user's machine and no prompt or screenshot is sent to a model provider}. It does not follow that no data leaves the machine, and we make no such claim. The desktop design addresses several of these gaps: a deterministic tier classifier instead of a model-assigned risk level, action-bound approvals, and a hash-chained audit log that covers local actions.

\paragraph{CAPTCHAs and automation etiquette} When the agent detects a CAPTCHA or bot-check page, it stops, marks the task \textsc{Blocked}, keeps the browser open, and offers the user two choices: solve the challenge and resume, or restart the search on an alternative engine. The resume path re-checks the page and refuses to continue while the challenge is still present. The code contains no solver. We note, however, that the browser is launched with automation-detection features disabled, a desktop user-agent string and the \texttt{navigator.webdriver} property hidden, and that direct navigation to search-result URLs is rewritten into typed queries. These choices make the agent resemble a human-driven browser; they are appropriate for a personal assistant acting on a user's own requests at human rate, but they should not be combined with high request volumes, and deployments should respect site terms of service. Chromium also runs without its process sandbox, which weakens isolation from malicious pages; enabling the sandbox is a straightforward hardening step.
